% Fourteen representative nodes; one physical template for all three timelines.
\begin{tikzpicture}[x=1mm,y=1mm,font=\figurefont\fontsize{8}{10}\selectfont,
 history card/.style={draw=FigureBlue!65,fill=white,rounded corners=2mm,line width=.65pt,minimum width=68mm,minimum height=18mm},
 history mark/.style={circle,minimum size=4mm,inner sep=0pt,draw=FigureBlue,fill=FigureBlueFill,line width=.7pt}]
\path[use as bounding box] (0,-5) rectangle (156,184);
\node[anchor=west,font=\figurefont\fontsize{8}{10}\selectfont,text=FigureInk] at (1,181) {概率图模型：分解、推断与生成};
\node[history card] at (34,169.0) {};
\node[anchor=west,font=\figurefont\fontsize{8}{10}\selectfont,text=FigureBlue] at (4,173.5) {1925\quad Ising模型};
\node[anchor=north west,align=left,text width=62mm,font=\FigureNoteFont\fontsize{7.5}{9.5}\selectfont,text=FigureInk] at (4,170.5) {局部相互作用定义能量\\归一化仍需汇总配置};
\node[history card] at (122,169.0) {};
\node[anchor=west,font=\figurefont\fontsize{8}{10}\selectfont,text=FigureBlue] at (92,173.5) {1953—1970\quad Metropolis—Hastings};
\node[anchor=north west,align=left,text width=62mm,font=\FigureNoteFont\fontsize{7.5}{9.5}\selectfont,text=FigureInk] at (92,170.5) {提议与接受组织状态更新\\绕过未知归一化常数};
\node[history card] at (34,147.5) {};
\node[anchor=west,font=\figurefont\fontsize{8}{10}\selectfont,text=FigureBlue] at (4,152.0) {1977\quad EM方法};
\node[anchor=north west,align=left,text width=62mm,font=\FigureNoteFont\fontsize{7.5}{9.5}\selectfont,text=FigureInk] at (4,149.0) {后验推断处理未观测量\\期望统计量参与参数更新};
\node[history card] at (122,147.5) {};
\node[anchor=west,font=\figurefont\fontsize{8}{10}\selectfont,text=FigureBlue] at (92,152.0) {1984\quad Gibbs图像恢复};
\node[anchor=north west,align=left,text width=62mm,font=\FigureNoteFont\fontsize{7.5}{9.5}\selectfont,text=FigureInk] at (92,149.0) {邻域关系描述像素依赖\\局部条件更新用于恢复};
\node[history card] at (34,126.0) {};
\node[anchor=west,font=\figurefont\fontsize{8}{10}\selectfont,text=FigureBlue] at (4,130.5) {1988\quad 贝叶斯网络};
\node[anchor=north west,align=left,text width=62mm,font=\FigureNoteFont\fontsize{7.5}{9.5}\selectfont,text=FigureInk] at (4,127.5) {条件独立赋予图语义\\有向分解组织联合分布};
\node[history card] at (122,126.0) {};
\node[anchor=west,font=\figurefont\fontsize{8}{10}\selectfont,text=FigureBlue] at (92,130.5) {1988\quad 团树局部计算};
\node[anchor=north west,align=left,text width=62mm,font=\FigureNoteFont\fontsize{7.5}{9.5}\selectfont,text=FigureInk] at (92,127.5) {图变换组织团与分隔集\\局部消息完成精确查询};
\node[history card] at (34,104.5) {};
\node[anchor=west,font=\figurefont\fontsize{8}{10}\selectfont,text=FigureBlue] at (4,109.0) {1996\quad 图模型理论};
\node[anchor=north west,align=left,text width=62mm,font=\FigureNoteFont\fontsize{7.5}{9.5}\selectfont,text=FigureInk] at (4,106.0) {系统联系分离与分解\\区分图结构与概率假设};
\node[history card] at (122,104.5) {};
\node[anchor=west,font=\figurefont\fontsize{8}{10}\selectfont,text=FigureBlue] at (92,109.0) {2001\quad 因子图};
\node[anchor=north west,align=left,text width=62mm,font=\FigureNoteFont\fontsize{7.5}{9.5}\selectfont,text=FigureInk] at (92,106.0) {连接变量与局部函数\\统一描述乘积与消息};
\node[history card] at (34,83.0) {};
\node[anchor=west,font=\figurefont\fontsize{8}{10}\selectfont,text=FigureBlue] at (4,87.5) {2008\quad 变分理论整合};
\node[anchor=north west,align=left,text width=62mm,font=\FigureNoteFont\fontsize{7.5}{9.5}\selectfont,text=FigureInk] at (4,84.5) {联系指数族与变分原理\\把近似推断写成优化};
\node[history card] at (122,83.0) {};
\node[anchor=west,font=\figurefont\fontsize{8}{10}\selectfont,text=FigureBlue] at (92,87.5) {2013—2014\quad 变分自编码器};
\node[anchor=north west,align=left,text width=62mm,font=\FigureNoteFont\fontsize{7.5}{9.5}\selectfont,text=FigureInk] at (92,84.5) {神经网络参数化生成分布\\学习摊销近似后验};
\node[history card] at (34,61.5) {};
\node[anchor=west,font=\figurefont\fontsize{8}{10}\selectfont,text=FigureBlue] at (4,66.0) {2015\quad 规范化流};
\node[anchor=north west,align=left,text width=62mm,font=\FigureNoteFont\fontsize{7.5}{9.5}\selectfont,text=FigureInk] at (4,63.0) {可逆映射改变密度\\Jacobian修正体积变化};
\node[history card] at (122,61.5) {};
\node[anchor=west,font=\figurefont\fontsize{8}{10}\selectfont,text=FigureBlue] at (92,66.0) {2015\quad 扩散生成模型};
\node[anchor=north west,align=left,text width=62mm,font=\FigureNoteFont\fontsize{7.5}{9.5}\selectfont,text=FigureInk] at (92,63.0) {逐步加噪定义前向链\\学习反向转移生成样本};
\node[history card] at (34,40.0) {};
\node[anchor=west,font=\figurefont\fontsize{8}{10}\selectfont,text=FigureBlue] at (4,44.5) {2020\quad 去噪扩散概率模型};
\node[anchor=north west,align=left,text width=62mm,font=\FigureNoteFont\fontsize{7.5}{9.5}\selectfont,text=FigureInk] at (4,41.5) {联系去噪与变分训练\\噪声预测支持逐步生成};
\node[history card] at (122,40.0) {};
\node[anchor=west,font=\figurefont\fontsize{8}{10}\selectfont,text=FigureBlue] at (92,44.5) {2022—2023\quad 流匹配};
\node[anchor=north west,align=left,text width=62mm,font=\FigureNoteFont\fontsize{7.5}{9.5}\selectfont,text=FigureInk] at (92,41.5) {回归条件向量场\\学习连续概率路径};
\begin{scope}[on background layer]
\draw[FigureMuted!35,line width=2.2mm,line cap=round] (73,169.0) -- (83,169.0) .. controls (91,164.0) and (65,152.5) .. (73,147.5) -- (83,147.5) .. controls (91,142.5) and (65,131.0) .. (73,126.0) -- (83,126.0) .. controls (91,121.0) and (65,109.5) .. (73,104.5) -- (83,104.5) .. controls (91,99.5) and (65,88.0) .. (73,83.0) -- (83,83.0) .. controls (91,78.0) and (65,66.5) .. (73,61.5) -- (83,61.5) .. controls (91,56.5) and (65,45.0) .. (73,40.0) -- (83,40.0);
\draw[white,dashed,line width=.3mm,dash pattern=on 1.5mm off 1.3mm] (73,169.0) -- (83,169.0) .. controls (91,164.0) and (65,152.5) .. (73,147.5) -- (83,147.5) .. controls (91,142.5) and (65,131.0) .. (73,126.0) -- (83,126.0) .. controls (91,121.0) and (65,109.5) .. (73,104.5) -- (83,104.5) .. controls (91,99.5) and (65,88.0) .. (73,83.0) -- (83,83.0) .. controls (91,78.0) and (65,66.5) .. (73,61.5) -- (83,61.5) .. controls (91,56.5) and (65,45.0) .. (73,40.0) -- (83,40.0);
\end{scope}
\node[history mark] at (73,169.0) {};
\draw[-{Stealth[length=1.4mm]},FigureStroke,line width=.6pt] (75.5,169.0) -- (80.5,169.0);
\node[history mark] at (83,169.0) {};
\node[history mark] at (73,147.5) {};
\draw[-{Stealth[length=1.4mm]},FigureStroke,line width=.6pt] (75.5,147.5) -- (80.5,147.5);
\node[history mark] at (83,147.5) {};
\node[history mark] at (73,126.0) {};
\draw[-{Stealth[length=1.4mm]},FigureStroke,line width=.6pt] (75.5,126.0) -- (80.5,126.0);
\node[history mark] at (83,126.0) {};
\node[history mark] at (73,104.5) {};
\draw[-{Stealth[length=1.4mm]},FigureStroke,line width=.6pt] (75.5,104.5) -- (80.5,104.5);
\node[history mark] at (83,104.5) {};
\node[history mark] at (73,83.0) {};
\draw[-{Stealth[length=1.4mm]},FigureStroke,line width=.6pt] (75.5,83.0) -- (80.5,83.0);
\node[history mark] at (83,83.0) {};
\node[history mark] at (73,61.5) {};
\draw[-{Stealth[length=1.4mm]},FigureStroke,line width=.6pt] (75.5,61.5) -- (80.5,61.5);
\node[history mark] at (83,61.5) {};
\node[history mark] at (73,40.0) {};
\draw[-{Stealth[length=1.4mm]},FigureStroke,line width=.6pt] (75.5,40.0) -- (80.5,40.0);
\node[history mark] at (83,40.0) {};
\node[anchor=west,font=\figurefont\fontsize{8}{10}\selectfont] at (1,27) {持续研究的问题};
\node[draw=FigureBlue,fill=FigureBlue!8!white,rounded corners=2mm,minimum width=49mm,minimum height=23mm,line width=.65pt] at (25,10) {};
\node[font=\figurefont\fontsize{8}{10}\selectfont] at (25,16) {结构与依赖};
\node[align=center,font=\FigureNoteFont\fontsize{7.5}{9.5}\selectfont] at (25,8) {结构学习、缺失观测\\识别条件};
\node[draw=FigureGreen,fill=FigureGreen!8!white,rounded corners=2mm,minimum width=49mm,minimum height=23mm,line width=.65pt] at (78,10) {};
\node[font=\figurefont\fontsize{8}{10}\selectfont] at (78,16) {推断与误差};
\node[align=center,font=\FigureNoteFont\fontsize{7.5}{9.5}\selectfont] at (78,8) {摊销近似、采样效率\\后验校准};
\node[draw=FigureYellow,fill=FigureYellow!8!white,rounded corners=2mm,minimum width=49mm,minimum height=23mm,line width=.65pt] at (131,10) {};
\node[font=\figurefont\fontsize{8}{10}\selectfont] at (131,16) {生成与控制};
\node[align=center,font=\FigureNoteFont\fontsize{7.5}{9.5}\selectfont] at (131,8) {条件分布、连续路径\\采样成本};
\end{tikzpicture}
